\documentclass[12pt]{article}
\usepackage[utf8]{inputenc}
\usepackage[T1]{fontenc}
\usepackage{hyperref}
\title{Ache Language Resources for K-12 Curriculum Development}
\author{Compiled by Claude Sonnet 5.5 (claude-sonnet-5-5) for Andr\'as Kornai}
\date{}
\begin{document}
\maketitle

\section{Introduction}
Ache (Ache Yi, also called Azhe or A-Che) is a Loloish language of south-central Yunnan, China, around Shuangbai, Yimen, Eshan and Lufeng counties. Its speakers are officially counted in the Yi nationality. Sources give figures of about 35,000 (Ethnologue, 2003) and 45,400 (2010). Search found only descriptive pages. It found no textbooks, primers, school programmes or audio.

\section{Primary Educational Resources}
None found.

\section{Linguistic and Reference Materials}
\begin{itemize}
\item Ache Yi language (Wikipedia): \url{https://en.wikipedia.org/wiki/Ache_Yi_language}
\item Glottolog entry (ache1244): \url{https://glottolog.org/resource/languoid/id/ache1244.bigmap.html}
\end{itemize}

\section{Religious and Community Resources}
\begin{itemize}
\item Asia Harvest people-group profile for the A-Che (a mission profile): \url{https://people-groups.asiaharvest.org/people-group-profiles/yunnan/a-che.pdf}
\end{itemize}

\section{Digital and Audio Resources}
None found.

\section{Curriculum Development Context}
No native-language educational material was found. Any curriculum would have to be built from scratch, probably with fieldwork. Standard Yi-script and Mandarin school materials are the likely regional context, but this search did not confirm them for Ache.
\end{document}